\section{Approach B: cluster, then inverted bitplanes and branching}
\label{sec:bitplanes}
Approach B keeps the same document score and common routing stage. Its extra index stores
the document bits by coordinate instead of only by row. That makes one sign test applicable
to many document positions at once.

\fig{bitplanes}{Approach B. A pending node has a set of active documents and an optimistic score bound.}{1}

\subsection{What a bitplane stores}
For the opened example cluster $\mathcal C_1=(2,5,7,8)$, the document codes are
\texttt{1011}, \texttt{1110}, \texttt{1001}, \texttt{0011}. The transposed index is
\begin{center}
\begin{tabular}{@{}lcccc@{}}\toprule
 & doc 2 & doc 5 & doc 7 & doc 8\\\midrule
bitplane 1 & 1&1&1&0\\
bitplane 2 & 0&1&0&0\\
bitplane 3 & 1&1&0&1\\
bitplane 4 & 1&0&1&1\\\bottomrule
\end{tabular}
\end{center}
Each row is a bitmap over document positions, not a list of full document vectors. An active
bitmap says which document positions still belong to the current node. To require coordinate
$i$ to be 1, intersect the active bitmap with plane $i$. To require 0, intersect with its
complement, while masking unused padding positions.

\begin{lstlisting}[caption={Build the additional bitplane index.}]
BuildBitplanes(cluster_codes, d, w):
    for each cluster j:
        W[j] = ceil(number_of_rows(j) / w)
        planes[j] = d arrays of W[j] zero words
        for local row u in j:
            for coordinate i in 0 .. d-1:
                if code[j, u, i] is 1:
                    set bit (u mod w) in planes[j, i, floor(u/w)]
    return planes
\end{lstlisting}

The extra persistent storage is exactly, for this layout,
\begin{equation}
M_{\rm planes}=d b_w\sum_{j=1}^{C}W_j,
\qquad M_{B,\rm index}=\Mbase+M_{\rm planes}.
\label{eq:planes-storage}
\end{equation}
Without per-cluster padding this is $Nd/8$, a second binary-code-sized payload.
For example, ten billion 256-bit codes occupy 320 GB; their bitplanes add approximately
another 320 GB, not 32 TB. Padding and other index structures are additional.

\subsection{Why the query suggests a traversal direction}
At coordinate $i$, choosing the query's sign adds no penalty; choosing the opposite sign adds
$a_i=|q_i|$. Sorting coordinates by descending $a_i$ tests the largest potential penalties first.
For the four-bit query, the order is $1,2,3,4$ with weights $0.7,0.4,0.2,0.1$.

In $\mathcal C_1$, matching bit 1 keeps documents $2,5,7$; the alternative contains document 8
with accumulated penalty $0.7$. Matching bit 2 next keeps $2,7$; document 5 moves to another
pending node with penalty $0.4$. This is a useful order, but neither alternative is automatically
safe to discard.

\begin{examplebox}{A greedy path can miss the best document}
Take $q=(0.6,0.5,0.4,0.3)$ and two documents: \texttt{0111} and \texttt{1000}.
The first has penalty $0.6$ and score $0.6$. The second has penalty $1.2$ and score $-0.6$.
A greedy choice of the matching first bit selects \texttt{1000} and misses the better document.
The largest individual weight is smaller than the sum of several remaining weights.
\end{examplebox}

\subsection{A bound makes branching safe}
Let a node have accumulated penalty $p$ on coordinates already tested. Untested coordinates
cannot contribute a negative mismatch penalty. Consequently,
\begin{equation}
U(\text{node})=A-2p
\label{eq:branch-bound}
\end{equation}
is an upper bound on every active document's final score. If the result heap is full and
$U$ is strictly below its worst retained score, that node cannot improve the result.
For equal scores, keep searching unless the ID tie rule is also certified.

The \textbf{frontier} is the collection of pending nodes. A best-first search takes the node
with the largest upper bound, or equivalently the smallest penalty. The bound is optimistic:
it may describe a sign combination that no actual document can achieve. A weak bound can keep
many nodes alive even though none ultimately improves the answer.

\begin{lstlisting}[caption={Dense bitplane search across the selected clusters.}]
SearchBitplanes(q, selected_clusters, eligible, table, R, L, B_node):
    order = coordinates sorted by decreasing abs(q[i])
    frontier = empty min-penalty priority queue
    best = empty top-R heap
    for each selected cluster j:
        active = bitmap of eligible rows, with physical width W[j]
        if active is not empty:
            frontier.push(j, depth=0, penalty=0, active)
    visits = 0
    while frontier is not empty:
        if best is full and best_possible(frontier) < best.worst_score:
            break
        if a finite B_node is set and visits == B_node:
            break
        node = frontier.pop_smallest_penalty()
        visits += 1
        if node.count <= L or node.depth == d:
            for word in node.active:             # includes zero words
                for each set bit in word:
                    score its document and offer it to best
            continue
        i = order[node.depth]
        matching = active AND preferred_sign_plane(j, i, q)
        opposite = active AND NOT preferred_sign_plane(j, i, q)
        queue nonempty matching child with the same penalty
        queue nonempty opposite child with penalty + abs(q[i])
    return best
\end{lstlisting}
$L$ is the leaf-size threshold and $B_{\rm node}$ the total node budget across all probes.
An unlimited budget permits exact termination by bounds; a finite budget can stop before the
remaining nodes are safely ruled out. A random exploration probability changes which pending
node is visited next. It is not a direct formula for recall or a synonym for the number of nodes.

\subsection{Count split visits, leaf visits, and scored documents separately}
For each selected cluster, let
\[
s_j=\text{visited split nodes},\quad
\ell_j=\text{visited leaf nodes},\quad
f_j=\text{documents fully scored}.
\]
Define $V=\sum_j(s_j+\ell_j)$ and $F=\sum_j f_j$. For a disjoint traversal without rescoring,
$F\leq\Nelig$. A finite node budget gives $V\leq B_{\rm node}$, but does not determine $F$.
The two principal bitmap counts are
\begin{equation}
U_{\rm split}=\sum_j s_j W_j,\qquad
U_{\rm leaf}=\sum_j \ell_j W_j.
\label{eq:bitmap-work}
\end{equation}
A split word in the reference implementation reads an active word and a plane word, produces
two child words, and counts matches. It includes two ANDs, a population count, complements and
an accumulation. Allocating zero-filled child arrays can add two initialization writes.
A leaf still scans all words in its active bitmap before enumerating set bits.

With separate effective coefficients,
\begin{align}
T_{B,\rm candidates}={}&T_{\rm order}+T_{\rm roots}
 +U_{\rm split}c_{\rm split}+U_{\rm leaf}c_{\rm leafword}
 +V c_{\rm node} \notag\\
 &+F G c_{\rm lut,leaf}+H_s(F,R)c_{\rm heap}+T_{\rm rank}(R),
\label{eq:branch-time}\\
\boxed{T_B=\Tcommon+T_{B,\rm candidates}+\Tfinish(R').}
\end{align}
Here $T_{\rm order}=O(d\log d)$ and $T_{\rm roots}=O(\sum_jW_j)$ for dense root construction.
$c_{\rm node}$ may be replaced by explicit queue-comparison, mask-allocation and copy terms.
Use one convention consistently; do not charge an allocation twice.

\subsection{The full-width cost is the important hidden term}
\begin{examplebox}{Few scores do not imply little work}
In the saved one-million-document run, $w=64$ gives $W=15{,}625$. A node budget of $32{,}768$
produced a mean $22{,}557.88$ split visits and $10{,}210.12$ leaf visits. Therefore
\[
U_{\rm split}=352{,}466{,}875,\quad U_{\rm leaf}=159{,}533{,}125.
\]
Their sum is 512 million word-loop visits. About $163{,}766$ documents were fully scored.
With $g=4$, that is 10.48 million score-table terms, versus 64 million terms for the scan.
The reduction in scoring was real, but it came with hundreds of millions of bitmap visits.
Word visits and score-table terms are not equal-cost instructions.
\end{examplebox}

If a tree were fully explored, perfectly balanced, and every internal node had two nonempty
children, then with $L$ documents per leaf there would be roughly $n_j/L$ leaves and one fewer
split. This gives a dense bitmap term of order
\begin{equation}
O\!\left(\sum_j\frac{n_j^2}{Lw}\right).
\label{eq:balanced-branch}
\end{equation}
This is an explanatory idealization, not a predictor for real embeddings. Fixed-bit splits can
have only one nonempty child. A finite traversal also leaves unvisited nodes. Thus the actual
split count need not be half the visited-node count.

A single greedy path visits only about $\log_2(n_j/L)$ levels in a balanced example, but still
pays a full $W_j$ at each level and has no general recall guarantee. In the broad worst case,
a fixed-coordinate trie can contain $O(n_jd)$ nonempty prefix nodes, making dense traversal
work as large as $O(d\sum_j n_j^2/w)$ before scoring and queue costs. Growing $N$ alone does
not guarantee that branching will overtake a scan.

\subsection{Peak query RAM: the pending masks matter}
Let $\mathcal F(t)$ be the frontier at time $t$, and let $j(v)$ identify a node's cluster.
If a node record needs $b_{\rm node}$ bytes and bitmap capacity is fully allocated,
\begin{equation}
Q_{B,\rm frontier}(t)=\sum_{v\in\mathcal F(t)}(b_{\rm node}+b_wW_{j(v)}).
\label{eq:frontier-memory}
\end{equation}
Capacity growth and allocator overhead can be represented by a factor $\gamma_B\geq1$ or,
preferably, by measured capacities. During a split, the parent and newly allocated children
can coexist. A conservative candidate-stage expression is
\begin{equation}
Q_{B,\rm search}=\Qbase+Q_{\rm order}
 +\max_t\left[\gamma_B Q_{B,\rm frontier}(t)+Q_{\rm current}(t)+Q_{\rm children}(t)\right].
\label{eq:branch-ram}
\end{equation}
For nonempty, disjoint pending sets, frontier count is at most $\Nelig$. Starting with at most
$P$ roots, each processed binary split can increase that count by at most one, so a finite
budget also gives $|\mathcal F|\leq P+B_{\rm node}$. Multiplying that bound by the largest
possible mask is conservative and can be very loose. It should not be reported as measured RAM.

\begin{resultbox}{When Approach B can help}
B can help when it avoids enough full scoring to repay the bitmap, queue and allocation work.
Useful conditions include strong early query penalties, small physical bitmaps, tight score
bounds and a modest live frontier. Unlimited search is exact inside the selected pool under
valid bounds and ties. Budget-limited search needs measured recall; its budget is not itself
an accuracy percentage.
\end{resultbox}
